\documentclass{situationist-preprint}

\title{The Bounds of Supererogation}
\msid{ETH-14-a}
\msfamily{ETH-14}
\msdiscipline{Ethics}
\mstype{Principal manuscript}
\mskeywords{supererogation; duty; decency; sainthood; heroism; agent-centred prerogative; self-sacrifice; demandingness}

\begin{document}
\maketitle

\begin{abstract}
An act is supererogatory if it is good to do, beyond what duty requires, and not wrong to omit. Some theories leave no room for such acts: maximizing consequentialism requires the best act, and a strict Kantian ethics requires every act that is good. We argue that there are supererogatory acts and that they occupy a bounded region. The lower bound is decency: the level of concern for others that duty requires, below which an omission is not a failure to be generous but a wrong. The upper bound is sainthood: the greatest sacrifice of the agent's own good that remains proportionate to the good it does for others and consistent with her other duties. Beyond it, further sacrifice ceases to be morally better and becomes waste or self-abasement. We derive the existence of supererogatory acts from an agent-centred prerogative, which lets agents give their own interests more than impartial weight in deciding whether to help, together with the distinction between reasons that require and reasons that justify. We then show that the prerogative alone fixes the lower bound and that a principle of proportion fixes the upper one.
\end{abstract}

\section{The problem}

A soldier throws himself on a grenade to save his comrades. A doctor leaves a comfortable practice to work in an epidemic. A neighbour spends her weekends caring for an elderly stranger. \citet{Urmson1958} argued that such acts are not explained by the traditional division of acts into the obligatory, the permitted and the forbidden. They are not obligatory, since we would not blame the soldier who took cover, and they are not merely permitted, since they are morally better than the alternatives and admirable in a way that permitted acts are not. Saints and heroes, in Urmson's terms, do more than duty requires.

Two kinds of theory deny that there is such a category. On maximizing consequentialism, an agent is required to do the best act available, so that no act is better than what duty requires; the soldier, if his act was best, was required to do it \citep{Kagan1989}. On some readings of Kant, imperfect duties of beneficence require us to adopt the end of others' happiness, and every act that serves this end without violating a perfect duty is a way of fulfilling a duty, so that nothing lies beyond it. Defenders of supererogation answered with accounts of why moral reasons for helping do not always require. \citet{Heyd1982} gave the classical analysis of the concept and grounded supererogation in the value of the agent's autonomy; \citet{Scheffler1982} argued for an agent-centred prerogative to give one's own projects more weight than their impartial value; \citet{Gert2003} and \citet{Portmore2011} distinguished the requiring strength of a reason from its justifying strength, and explained supererogation as the case in which the agent's own interests justify omission without requiring it.

These accounts explain why there is a region above duty. They say little about its shape. In particular, they do not say whether the region is bounded above: whether every additional sacrifice for others is more admirable than the last. \citet{Wolf1982} argued that the life of a moral saint, devoted at every point to the good of others, is not an ideal, and \citet{Hill1973} that a self-abnegation which gives up one's own standing is a vice. We argue that these observations mark a boundary of supererogation, not merely a limit on how attractive saints are, and that supererogation is bounded above by sainthood and below by decency.

\section{Framework}

Consider an agent in a situation with a set of available acts. For each act $a$, let $\sigma(a)$ be the sacrifice it imposes on the agent, measured in her own well-being relative to the least costly available act, and let $g(a)$ be the good it does for others, on the same scale.

\begin{definition}[Supererogation]\label{d:super}
An act is \emph{supererogatory} if it is permissible, its omission is permissible, and it is morally better than some permissible alternative \citep{Heyd1982}.
\end{definition}

\begin{principle}[Prerogative]\label{p:prerog}
In deciding whether to benefit others, an agent may weigh her own sacrifice at more than its impartial value, by a factor that grows with the size of the sacrifice. She is required to benefit others when the good done is great relative to the sacrifice so weighted, and is not required otherwise.
\end{principle}

Principle~\ref{p:prerog} is Scheffler's prerogative in the form needed here \citep{Scheffler1982}. Its consequence is a threshold of required sacrifice that depends on the stakes. Small sacrifices for great goods are required; large sacrifices are not, unless the good is overwhelming. The rescue of a drowning child at the cost of a ruined suit is required; the rescue of the same child at a serious risk to one's life is not. The results of ETH-18 (a duty to rescue up to moderate cost) and ETH-06 (a duty to give a substantial fixed share of income, but not to the point of marginal utility) are instances of the principle.

\begin{definition}[Decency]\label{d:decency}
The \emph{decency threshold} in a situation is the level of sacrifice that duty requires of the agent in it, by Principle~\ref{p:prerog}. An act is \emph{decent} if it does at least what duty requires.
\end{definition}

The term is from \citet{Thomson1971}, who distinguished the Good Samaritan from the Minimally Decent Samaritan: the first went out of his way at some cost; the second does what decency requires and no more, and failing to do even that is a wrong.

\begin{principle}[Requiring and justifying]\label{p:justify}
The prerogative justifies omitting a beneficial act without making the act worse: an act that does more good for others at greater cost to the agent remains morally better than its alternatives, though not required, so long as the good it does is not outweighed by the sacrifice it imposes.
\end{principle}

\begin{principle}[Proportion]\label{p:prop}
A sacrifice is morally good only if the good it does for others is at least as great as the sacrifice, assessed impartially, and only if it does not violate the agent's other duties, including duties to those who depend on her and to herself.
\end{principle}

Principle~\ref{p:prop} says that self-sacrifice is valuable for what it achieves, not in itself. A sacrifice that achieves less than it costs, impartially assessed, makes the world worse; a sacrifice that abandons one's children or one's own standing as a person is not made good by the benefit it confers on others.

\begin{definition}[Sainthood]\label{d:saint}
The \emph{saintly bound} in a situation is the greatest sacrifice that satisfies Principle~\ref{p:prop}.
\end{definition}

The saint, on this definition, is not the person who sacrifices most, but the person who sacrifices as much as is good to sacrifice. The definition fits the saints of the moral tradition better than the alternative: what is admired in them is the good they did at great cost, not the cost.

\section{The argument}

\begin{lemma}[Lower bound]\label{l:lower}
No act whose sacrifice is at or below the decency threshold is supererogatory.
\end{lemma}

\begin{proof}
By Definition~\ref{d:decency}, duty requires the agent to make sacrifices up to the threshold when the corresponding good is at stake. An act that does less is impermissible. An act that does only what duty requires is not one whose omission is permissible, since omitting it in favour of doing less violates the duty; so it is not supererogatory by Definition~\ref{d:super}.
\end{proof}

\begin{lemma}[Upper bound]\label{l:upper}
No act whose sacrifice exceeds the saintly bound is supererogatory.
\end{lemma}

\begin{proof}
By Definition~\ref{d:saint}, an act whose sacrifice exceeds the saintly bound violates Principle~\ref{p:prop}: either the good it does for others is less than the sacrifice, impartially assessed, or it violates one of the agent's other duties. In the first case it is not morally better than the alternative of a smaller sacrifice, since the extra sacrifice costs more than it achieves; in the second it is impermissible. In neither case does it meet Definition~\ref{d:super}.
\end{proof}

\begin{theorem}\label{t:main}
There are supererogatory acts. An act is supererogatory exactly when it is permissible, satisfies Principle~\ref{p:prop}, and imposes a sacrifice above the decency threshold; every such act imposes a sacrifice at or below the saintly bound.
\end{theorem}

\begin{proof}
By Lemma~\ref{l:lower}, a supererogatory act imposes a sacrifice above the decency threshold. It satisfies Principle~\ref{p:prop}, since otherwise, as in the proof of Lemma~\ref{l:upper}, it would be either impermissible or no better than a smaller sacrifice; and by Definition~\ref{d:saint} an act satisfying the principle imposes a sacrifice at or below the saintly bound. Conversely, let $a$ be a permissible act that satisfies Principle~\ref{p:prop} and whose sacrifice lies above the decency threshold. Its omission is permissible: by Principle~\ref{p:prerog} the agent is not required to sacrifice beyond the threshold, and doing only what decency requires is permissible. Since the good $a$ does is at least its sacrifice, Principle~\ref{p:justify} makes it morally better than the merely decent alternative. By Definition~\ref{d:super}, $a$ is supererogatory. Such acts exist whenever the saintly bound lies above the decency threshold, as it does in the cases of the soldier, the doctor and the neighbour: each makes a sacrifice beyond what duty requires, for a good to others that is at least as great.
\end{proof}

The theorem gives supererogation a definite place in the structure of morality. Below the decency threshold lie the wrongs of omission: refusing a costless rescue, keeping everything while others starve. Between the thresholds lie the acts of generosity and heroism that are good to do and not wrong to omit; ETH-03-a counts self-sacrifice to save five among them. Above the saintly bound lie sacrifices that are no longer good: the person who gives away what her children need, who risks her life to save a stranger's property, or who effaces herself until she no longer counts her own interests at all. The demandingness of morality, which ETH-28 argues is moderate, is set by the lower bound; its aspirations, by the region above it.

\section{Objections and replies}

\begin{objection}[The paradox of supererogation]\label{o:paradox}
If a supererogatory act is morally better than the alternatives, there is a stronger moral reason to do it, and an agent who acts against the stronger moral reason acts wrongly. So supererogatory acts are required after all \citep{Kagan1989}.
\end{objection}

\begin{reply}
The objection assumes that the only reasons relevant to whether an act is required are moral reasons of the kind that make acts better. Principle~\ref{p:justify} denies this. The agent's own interests, weighted as Principle~\ref{p:prerog} permits, have justifying force: they make it permissible not to do the better act, without making the better act worse. That reasons can justify without requiring is shown by the structure of permission in other domains: one may take the harder route up a mountain for the view, though the easier route is permissible \citep{Gert2003,Portmore2011}.
\end{reply}

\begin{objection}[The more the better]\label{o:more}
Admiration for self-sacrifice increases with its extent. The soldier who throws himself on the grenade is admired more than one who risks a wound. There is no upper bound: the greater the sacrifice, the greater the moral worth.
\end{objection}

\begin{reply}
The soldier's sacrifice is within the saintly bound: he gives his life for several lives. The test of the objection is a case of disproportion, and there admiration fails. The man who drowns trying to retrieve a stranger's hat, or the mother who gives away her children's food to strangers no needier than they, are not admired more for sacrificing more. What is admired in sacrifice is the good it secures together with the cost, and when the cost exceeds the good the sacrifice is pitied or regretted. This is what Lemma~\ref{l:upper} records. \citet{Wolf1982} and \citet{Hill1973} describe two ways in which the bound is overstepped: by neglect of one's own goods and those of one's intimates, and by servility, the failure to count one's own interests as interests of a person with equal standing.
\end{reply}

\begin{objection}[Small favours]\label{o:favours}
Many supererogatory acts involve no sacrifice: holding a door, lending a book, giving directions at a slight detour. If the lower bound is a threshold of sacrifice, these acts fall below it and should be required, yet they are not.
\end{objection}

\begin{reply}
The decency threshold is not a fixed quantity of sacrifice; it is the level of sacrifice that duty requires in the situation, and by Principle~\ref{p:prerog} that level depends on the good at stake. Where the good is small, duty may require no sacrifice at all, and the threshold is zero. A small favour then lies above it, and is supererogatory in a minor way, as \citet{Heyd1982} holds. The structure of Theorem~\ref{t:main} is unchanged: the favour is permissible, its omission is permissible, and it does more good than it costs.
\end{reply}

\begin{objection}[Kantian duties of beneficence]\label{o:kant}
The duty of beneficence is imperfect: it requires adopting the happiness of others as an end, and leaves latitude in how and when to pursue it. Acts above the threshold are ways of fulfilling this duty, not acts beyond it.
\end{objection}

\begin{reply}
The latitude of an imperfect duty is exactly what makes room for supererogation. If a given act of help is one of many ways of fulfilling the duty, and the agent may permissibly choose another, then the act is not required, and if it does more good at greater cost than the minimum that fulfils the duty, it is better than a permissible alternative. That it also counts as fulfilling the duty does not prevent it from meeting Definition~\ref{d:super}, which requires only that its omission be permissible.
\end{reply}

\section{What the result does not establish}

Theorem~\ref{t:main} is structural. It does not fix where the decency threshold lies in particular cases, which is the subject of ETH-06 and ETH-18, or how the prerogative's weighting factor grows with the size of the sacrifice. It does not say that every act within the bounds is equally admirable; within the region, acts that do more good at greater cost are generally better. Nor does it say that the saintly bound is easy to locate: in particular cases it may be unclear whether a sacrifice is proportionate. The theorem says only that there is such a bound and that acts beyond it are not supererogatory.

\section{Conclusion}

There are supererogatory acts. They lie above decency, the level of sacrifice that duty requires given an agent-centred prerogative, and at or below sainthood, the greatest sacrifice that remains proportionate to the good it does and consistent with the agent's other duties. Below decency lie wrongs of omission; above sainthood, sacrifices that are no longer good. The region between them is where the heroic and the saintly acts of the moral tradition belong.

\begin{thebibliography}{99}\small
\bibitem[Gert(2003)]{Gert2003} Gert, J. (2003). Requiring and justifying: Two dimensions of normative strength. \emph{Erkenntnis}, 59(1), 5--36.
\bibitem[Heyd(1982)]{Heyd1982} Heyd, D. (1982). \emph{Supererogation: Its Status in Ethical Theory}. Cambridge University Press.
\bibitem[Hill(1973)]{Hill1973} Hill, T.~E., Jr. (1973). Servility and self-respect. \emph{The Monist}, 57(1), 87--104.
\bibitem[Kagan(1989)]{Kagan1989} Kagan, S. (1989). \emph{The Limits of Morality}. Clarendon Press.
\bibitem[Portmore(2011)]{Portmore2011} Portmore, D.~W. (2011). \emph{Commonsense Consequentialism: Wherein Morality Meets Rationality}. Oxford University Press.
\bibitem[Scheffler(1982)]{Scheffler1982} Scheffler, S. (1982). \emph{The Rejection of Consequentialism}. Clarendon Press.
\bibitem[Thomson(1971)]{Thomson1971} Thomson, J.~J. (1971). A defense of abortion. \emph{Philosophy \& Public Affairs}, 1(1), 47--66.
\bibitem[Urmson(1958)]{Urmson1958} Urmson, J.~O. (1958). Saints and heroes. In A.~I. Melden (Ed.), \emph{Essays in Moral Philosophy}. University of Washington Press.
\bibitem[Wolf(1982)]{Wolf1982} Wolf, S. (1982). Moral saints. \emph{The Journal of Philosophy}, 79(8), 419--439.
\end{thebibliography}

\end{document}
